% Exported from Math++. Compile with pdflatex.
\documentclass[11pt,a4paper]{article}
\usepackage[T1]{fontenc}
\usepackage[utf8]{inputenc}
\usepackage{lmodern}
\usepackage{amsmath,amssymb,amsthm,mathtools}
\usepackage{graphicx,booktabs,xcolor}
\usepackage[hidelinks]{hyperref}
\usepackage[margin=25mm]{geometry}
\providecommand{\degree}{\ensuremath{^\circ}}
\providecommand{\xmark}{\ensuremath{\times}}
\providecommand{\omicron}{o}
\providecommand{\oiint}{\mathop{\oint\mkern-12mu\oint}}
\theoremstyle{plain}
\newtheorem{theorem}{Theorem}
\newtheorem{lemma}[theorem]{Lemma}
\newtheorem{corollary}[theorem]{Corollary}
\newtheorem{proposition}[theorem]{Proposition}
\newtheorem{conjecture}[theorem]{Conjecture}
\theoremstyle{definition}
\newtheorem{definition}[theorem]{Definition}
\newtheorem{example}[theorem]{Example}
\theoremstyle{remark}
\newtheorem*{remark}{Remark}
\setlength{\parindent}{0pt}\setlength{\parskip}{0.6em}
\title{Numerical methods}
\date{}
\begin{document}
\maketitle
Each method shows its answer with an error estimate; \texttt{= ?!} shows every step.

\section*{Riemann sums, trapezoid and Simpson}

\begin{equation*}
  \operatorname{Riemann\ sum}_{left}(x^2, [0, 1], 4) = \frac{7}{32}
\end{equation*}
\begin{equation*}
  \operatorname{Riemann\ sum}_{right}(x^2, [0, 1], 4) = \frac{15}{32}
\end{equation*}
\begin{equation*}
  \operatorname{Midpoint\ rule}(x^2, [0, 1], 4) = \frac{21}{64}
\end{equation*}
\begin{equation*}
  \operatorname{Trapezoid\ rule}(\sin x, [0, \pi], 10) \approx 1.983523538
\end{equation*}
\begin{equation*}
  \operatorname{Simpson's\ rule}(e^{-x^2}, [0, 1], 10) \approx 0.7468249483
\end{equation*}
\begin{equation*}
  \operatorname{Simpson's\ rule}(x^3, [0, 2], 2) = 4
\end{equation*}
\section*{Newton's method}

\begin{equation*}
  \operatorname{Newton's\ method}(x^2 - 2, 1) \approx 1.414213562
\end{equation*}
\begin{equation*}
  \operatorname{Newton's\ method}(\cos x - x, 1) \approx 0.7390851332
\end{equation*}
\section*{Euler's method and RK4}

For $y' = y$ with $y(0) = 1$, step $h = 0.1$ and ten steps, the exact value is $y(1) = e$.

\begin{equation*}
  \operatorname{Euler's\ method}(y, 0, 1, 0.1, 10) \approx 2.59374246
\end{equation*}
\begin{equation*}
  \operatorname{RK4}(y, 0, 1, 0.1, 10) \approx 2.718279744
\end{equation*}
\begin{equation*}
  \operatorname{RK4}(x + y, 0, 1, 0.1, 10) \approx 3.436559488
\end{equation*}
\section*{Taylor error bounds}

The Lagrange bound for the cubic Taylor polynomial of $e^x$ about $0$ at $x = 1$, and the degree needed for $\sin x$ at $x = 0.5$ to be within $10^{-6}$.

\begin{equation*}
  \operatorname{Taylor\ bound}(e^x, 0, 3, 1) \approx 0.1132617429
\end{equation*}
\begin{equation*}
  \operatorname{Terms\ needed}(\sin x, 0, 0.5, 0.000001) = 7
\end{equation*}
\section*{Convergence tests}

Every test is reported by name.

\begin{equation*}
  \operatorname{tests}(\frac{1}{k^2}) = \text{true}
\end{equation*}
\begin{equation*}
  \operatorname{tests}(\frac{(-1)^k}{k}) = \text{true}
\end{equation*}
\begin{equation*}
  \operatorname{tests}(\frac{2^k}{k!}) = \text{true}
\end{equation*}
\section*{Radius and interval of convergence}

\begin{equation*}
  \operatorname{ioc}(\frac{(x-2)^k}{k 3^k}) = \left[-1, 5\right)
\end{equation*}
\begin{equation*}
  \operatorname{ioc}(\frac{x^k}{k^2}) = \left[-1, 1\right]
\end{equation*}
\begin{equation*}
  \operatorname{Radius\ of\ convergence}(\frac{x^k}{k!}) = \infty
\end{equation*}
\end{document}
